% ============================================================================
% Fichier  : partie3/P03_C08_duplication_hachage.tex : FORN (LPI342)
% Sources : ISO/IEC 27037 ; NIST SP 800-86 ; ACPO. Opération TSANAGA.
% ============================================================================

\chapter{Duplication et Hachage}
\label{chap:duplication}
\chaptermark{Duplication et hachage}

\epigraph{%
    « Le condensat est la signature de l'intégrité~:\\
    un bit change, il change tout. »%
}{Principe du cours}

\niveauChapitre{Intermédiaire}{Chapitre~\ref{chap:preserver}.}

\begin{boxobjectifs}
\begin{itemize}
    \item Comprendre la copie bit à bit (image).
    \item Comprendre les fonctions de hachage (MD5, SHA-256).
    \item Utiliser le condensat pour prouver l'intégrité.
    \item Travailler sur la copie, jamais sur l'original.
    \item Documenter l'acquisition.
\end{itemize}
\end{boxobjectifs}

\bigskip

% ============================================================================
\section*{Question directrice : comment prouver que la copie du disque est identique à l'original, et le restera ?}
% ============================================================================

On n'analyse jamais l'original~: on en fait une image, dont on prouve la
fidélité par un condensat. Ce chapitre explique la duplication et le hachage,
cœur technique de l'intégrité.

\medskip
\textbf{Plan du chapitre.}
\begin{enumerate}
    \item La copie bit à bit (section~\ref{sec:c08-image}).
    \item Les fonctions de hachage (section~\ref{sec:c08-hachage}).
    \item Prouver l'intégrité (section~\ref{sec:c08-integrite}).
    \item Documenter l'acquisition (section~\ref{sec:c08-documenter}).
\end{enumerate}

\separateur

% ============================================================================
\section{La Copie Bit à Bit}
\label{sec:c08-image}
% ============================================================================

\begin{boxdef}
\termecle{Image forensique}~: copie \emph{bit à bit} d'un support, qui
reproduit non seulement les fichiers visibles, mais aussi l'espace non alloué,
les fragments et les métadonnées. Elle se distingue d'une simple copie de
fichiers, qui perd l'essentiel des traces.
\end{boxdef}

% ============================================================================
\section{Les Fonctions de Hachage}
\label{sec:c08-hachage}
% ============================================================================

\begin{boxdef}
\termecle{Fonction de hachage}~: fonction qui transforme une donnée, quelle
que soit sa taille, en une empreinte de longueur fixe (le \termecle{condensat}).
Un changement d'un seul bit produit un condensat entièrement différent. MD5 et
SHA-256 en sont des exemples~; SHA-256 est préféré, MD5 étant considéré comme
faible pour les usages sensibles.
\end{boxdef}

\begin{tableaucours}{L{3.0cm}L{9.2cm}}{%
    \tableauHeader{Fonction} &
    \tableauHeader{Usage}
}
MD5 & Historique, rapide~; éviter seule pour un usage sensible \\
SHA-256 & Recommandée pour l'intégrité forensique \\
\end{tableaucours}
\vspace{-0.8em}
\begin{center}{\small\itshape Tableau~8.1. Deux fonctions de hachage}\end{center}

% ============================================================================
\section{Prouver l'Intégrité}
\label{sec:c08-integrite}
% ============================================================================

On calcule le condensat de l'original \emph{avant} la copie, puis celui de la
copie. S'ils sont identiques, la copie est fidèle et l'original n'a pas
changé. On recalcule le condensat de la copie avant chaque usage~: s'il
correspond toujours, la copie est restée intègre.

\begin{boxCRO}
\textbf{Le condensat sert l'opposabilité.} Un condensat calculé et consigné
établit la \cro{0}{+}{+}~: il prouve l'intégrité (fiabilité) et rend la pièce
défendable (opposabilité). Sans condensat, l'intégrité n'est qu'affirmée.
\end{boxCRO}

% ============================================================================
\section{Documenter l'Acquisition}
\label{sec:c08-documenter}
% ============================================================================

\begin{boxscenario}{Opération TSANAGA. L'image du disque}
\textbf{Constatation.} Derrière un bloqueur en écriture, Mme~NKENG réalise
une image bit à bit du disque de l'ordinateur et calcule son condensat
SHA-256, qu'elle consigne. Toute l'analyse se fera sur cette image.

\medskip
\textbf{Analyse} (confiance élevée). Le condensat consigné permettra, à tout
moment, de prouver que l'image analysée est fidèle à l'original. L'original
est remis sous scellé, intact.
\end{boxscenario}

\begin{boxobjet}{Image disque TSA-2026-01}
\textbf{Nature~:} image bit à bit du disque de l'ordinateur saisi.
\textbf{Source~:} disque TSA-2026-01, derrière bloqueur en écriture.
\textbf{Intégrité~:} condensat SHA-256 consigné à l'acquisition.
\textbf{Opposabilité~:} chaîne de possession continue. \textbf{Usage~:}
analyse sur l'image, original scellé.
\end{boxobjet}

\begin{boxmetier}{Analyste forensique}
\textbf{Hacher, copier, re-hacher.} L'analyste calcule le condensat avant et
après la copie, et le recalcule avant chaque session de travail. Trois
gestes, une garantie d'intégrité à toute épreuve.
\end{boxmetier}

\separateur

\begin{boxresume}
\textbf{Ce qu'il faut retenir de ce chapitre~:}
\begin{itemize}
    \item \textbf{L'image bit à bit} copie tout, y compris l'espace non
          alloué.
    \item \textbf{Le hachage} produit un condensat sensible au moindre bit.
    \item \textbf{SHA-256} est préféré à MD5 pour l'intégrité.
    \item \textbf{Le condensat} prouve la fidélité de la copie.
    \item \textbf{On analyse la copie}, l'original reste scellé.
\end{itemize}
\end{boxresume}

\begin{boxinfo}
\textbf{Points de contrôle de l'élève.} Je sais~: définir l'image bit à bit~;
expliquer le hachage~; prouver l'intégrité par condensat~; documenter une
acquisition.
\end{boxinfo}

\bigskip

\begin{boxexo}[title={\ding{45}\quad Exercice~8.1. Image \niveaubase}]
Qu'est-ce qu'une image bit à bit, et en quoi diffère-t-elle d'une copie de
fichiers~?
\end{boxexo}

\begin{boxexo}[title={\ding{45}\quad Exercice~8.2. Condensat \niveaubase}]
Que se passe-t-il pour le condensat si un seul bit change~?
\end{boxexo}

\begin{boxexo}[title={\ding{45}\quad Exercice~8.3. MD5 ou SHA-256 \niveaubase}]
Quelle fonction préférer pour l'intégrité forensique~? Pourquoi~?
\end{boxexo}

\begin{boxexo}[title={\ding{45}\quad Exercice~8.4. Prouver \niveaumoyen}]
Comment prouver que la copie est fidèle et le reste~?
\end{boxexo}

\begin{boxexo}[title={\ding{45}\quad Exercice~8.5. Avant la copie \niveaumoyen}]
Pourquoi calculer le condensat de l'original avant la copie~?
\end{boxexo}

\begin{boxexo}[title={\ding{45}\quad Exercice~8.6. L'espace non alloué \niveaumoyen}]
Pourquoi une image bit à bit capture-t-elle l'espace non alloué~?
\end{boxexo}

\begin{boxexo}[title={\ding{45}\quad Exercice~8.7. CRO \niveaumoyen}]
Expliquez l'analyse CRO du condensat.
\end{boxexo}

\begin{boxexo}[title={\ding{45}\quad Exercice~8.8. Re-hacher \niveaumoyen}]
Pourquoi recalculer le condensat avant chaque session de travail~?
\end{boxexo}

\begin{boxexo}[title={\ding{45}\quad Exercice~8.9. Procédure \niveauexpert}]
Rédigez la procédure d'acquisition par image du disque de TSANAGA.
\end{boxexo}

\begin{boxexo}[title={\ding{45}\quad Exercice~8.10. Fiche d'image \niveauexpert}]
Concevez la fiche d'acquisition (image) du cabinet NYONG SÉCURITÉ.
\end{boxexo}

\bigskip

\begin{boxinfo}
\textbf{Pour aller plus loin.}
\begin{itemize}
    \item ISO/IEC 27037~\cite{iso27037}~; NIST SP~800-86~\cite{nist80086}.
\end{itemize}
\end{boxinfo}

\begin{boxinfo}
\textbf{Transition.} Toutes les traces ne survivent pas également~: le
Chapitre~\ref{chap:acquisition} traite de l'ordre de volatilité et de
l'acquisition à chaud.
\end{boxinfo}
